\section*{Chapter \ref{ch:b1:stereochemistry-in-depth} --- Stereochemistry in Depth}

\begin{solution}{exo:b1:stereochemistry-in-depth:1}
\omterm{def:b1:stereochemistry-in-depth:isomers}{Constitutional isomers}; \omterm{def:b1:stereochemistry-in-depth:enantiomers}{enantiomers}; \omterm{def:b1:stereochemistry-in-depth:enantiomers}{diastereomers} (\omterm{def:b1:stereochemistry-in-depth:isomers}{stereoisomers} that are
not mirror images); two \omterm{def:b1:stereochemistry-in-depth:isomers}{conformations} of one compound; \omterm{def:b1:stereochemistry-in-depth:enantiomers}{diastereomers}
(\cip{R,S} is the meso form).
\end{solution}

\begin{solution}{exo:b1:stereochemistry-in-depth:2}
\ce{-OH} $>$ \ce{-NH2} $>$ \ce{-CH3} $>$ \ce{-H} (O, N, C, H).
\ce{-COOH} (O,O,O) $>$ \ce{-CHO} (O,O,H) $>$ \ce{-CH2OH} (O,H,H) $>$
\ce{-CH(CH3)2} (C,C,H). \ce{-Br} $>$ \ce{-Cl} $>$ \ce{-CCl3} (Cl,Cl,Cl)
$>$ \ce{-CH2Cl} (Cl,H,H): the first atom decides before its neighbours.
\end{solution}

\begin{solution}{exo:b1:stereochemistry-in-depth:3}
2,3-Dichlorobutane: two equivalent centres, \cip{R,R}, \cip{S,S} and the
meso \cip{R,S}: three. 2-Bromo-3-chlorobutane: two different centres,
four \omterm{def:b1:stereochemistry-in-depth:isomers}{stereoisomers}, no meso form. Pentane-2,4-diol: three, \cip{2R,4S}
being meso.
\end{solution}

\begin{solution}{exo:b1:stereochemistry-in-depth:4}
On C5: isopropenyl (C,C,C, the double bond counting twice) $>$ C6 $>$ C4
$>$ H. C6 and C4 both carry (C,H,H); one step further C6 leads to the
carbonyl carbon (O,O,C) and C4 to C3 (C,C,H): C6 wins. With the isopropenyl
group in front and H behind, the order isopropenyl $\to$ C6 $\to$ C4
turns clockwise as drawn: \cip{R}, spearmint carvone; the hashed one is
\cip{S}.
\end{solution}

\begin{solution}{exo:b1:stereochemistry-in-depth:5}
$\theta = 0$: methyls eclipsed, \qty{16.5}{kJ/mol}; $60^\circ$: gauche, about
\qty{2.9}{kJ/mol} (minimum 2.8 at $62^\circ$); $120^\circ$: methyl eclipsing H,
\qty{15.2}{kJ/mol}; $180^\circ$: anti, 0. With equal energies, two \omterm{def:b1:stereochemistry-in-depth:conformations}{gauche
conformations} for one anti: two thirds gauche.
\end{solution}

\begin{solution}{exo:b1:stereochemistry-in-depth:6}
Chain vertical with \ce{COOH} on top and \ce{CH3} at the bottom; ranks
$\ce{NH2} > \ce{COOH} > \ce{CH3} > \ce{H}$. Putting \ce{NH2} on the left
and H on the right, $\ce{NH2} \to \ce{COOH} \to \ce{CH3}$ turns clockwise
as seen; H is horizontal (towards the reader), so the descriptor is
reversed: \cip{S}. The amino group is on the left.
\end{solution}

\begin{solution}{exo:b1:stereochemistry-in-depth:7}
\textit{cis}: on adjacent carbons one bond up and one down among the \omterm{def:b1:stereochemistry-in-depth:conformations}{axial}
and equatorial positions gives axial--equatorial in both \omterm{def:b1:stereochemistry-in-depth:conformations}{chairs}.
\textit{trans}: diequatorial or diaxial. In the diequatorial \omterm{def:b1:stereochemistry-in-depth:conformations}{chair} the two
methyls are gauche to each other (one interaction); the \textit{cis}
isomer has that interaction plus an \omterm{def:b1:stereochemistry-in-depth:conformations}{axial} methyl, about \qty{5.7}{kJ/mol}
more. The \textit{trans} isomer is the more stable.
\end{solution}

\begin{solution}{exo:b1:stereochemistry-in-depth:8}
$[\alpha] = 13.2/(2.00 \times 0.100) = +66.0^\circ$. A \qty{1.00}{dm} tube:
$+6.60^\circ$. The enantiomer: $-13.2^\circ$ in the \qty{2.00}{dm} tube.
\end{solution}

\begin{solution}{exo:b1:stereochemistry-in-depth:9}
\qty{25}{\celsius}: $\exp(5700/(8.314 \times 298.15)) = 10.0$, fraction
$10.0/11.0 = \qty{91}{\percent}$. \qty{-50}{\celsius}: $\exp(5700/(8.314
\times 223.15)) = 21.6$, \qty{96}{\percent}.
\end{solution}

\begin{solution}{exo:b1:stereochemistry-in-depth:10}
In ethane all six hydrogens are alike: every \omterm{def:b1:stereochemistry-in-depth:conformations}{staggered conformation} is
the same. In butane the rear methyl can stand opposite the front one
(anti) or beside it ($\pm 60^\circ$, gauche), where the two methyls crowd
each other. Populations: anti 1, gauche $2\exp(-2.8/2.48) = 0.65$: about
\qty{61}{\percent} anti and \qty{39}{\percent} gauche.
\end{solution}

\begin{solution}{exo:b1:stereochemistry-in-depth:11}
C2 and C4 are stereogenic. C3 carries H, OH and two branches that are
identical in constitution; they differ only when C2 and C4 have opposite
descriptors, and only then is C3 stereogenic (pseudo-asymmetric). Isomers:
\cip{2R,4R} and \cip{2S,4S}, a pair of \omterm{def:b1:stereochemistry-in-depth:enantiomers}{enantiomers} (C3 not stereogenic);
\cip{2R,4S} with the two arrangements at C3, two meso forms. Four in all.
\end{solution}

\begin{solution}{exo:b1:stereochemistry-in-depth:12}
$[\alpha]_{\text{sample}} = 2.31/(1.00 \times 0.050) = +46.2^\circ$, and
$2x - 1 = 46.2/66.0 = 0.700$: $x = 0.850$, \qty{85}{\percent} of the
$(+)$ enantiomer and \qty{15}{\percent} of the $(-)$.
\end{solution}

\begin{solution}{pb:b1:stereochemistry-in-depth:1}
\textbf{1.} Cyclohexane ring: C1 carries OH, C2 the isopropyl group, C5 the
methyl; C1, C2 and C5 are stereogenic.
\textbf{2.} $2^3 = 8$, four pairs of \omterm{def:b1:stereochemistry-in-depth:enantiomers}{enantiomers}.
\textbf{3.} No arrangement has a mirror plane: the three substituents are
all different.
\textbf{4.} \ce{OH} $>$ C2 (C,C,H) $>$ C6 (C,H,H) $>$ H.
\textbf{5.} C1 (O,C,H) $>$ isopropyl (C,C,H) $>$ C3 (C,H,H) $>$ H.
\textbf{6.} C6 and C4 both (C,H,H); next, C6 leads to C1 (O,C,H), C4 to C3
(C,H,H): C6 $>$ C4 $>$ \ce{CH3} $>$ H.
\textbf{7.} \cip{1S,2R,5S}.
\textbf{8.} In a \omterm{def:b1:stereochemistry-in-depth:conformations}{chair}, two \omterm{def:b1:stereochemistry-in-depth:conformations}{equatorial} bonds on adjacent carbons (C1, C2)
point one up and one down: \textit{trans}; on carbons 1 and 3 of the ring
(C1 and C5, through C6) both \omterm{def:b1:stereochemistry-in-depth:conformations}{equatorial} bonds point the same way:
\textit{cis}. That is the arrangement of menthol.
\textbf{9.} All three become \omterm{def:b1:stereochemistry-in-depth:conformations}{axial}.
\textbf{10.} In the flipped \omterm{def:b1:stereochemistry-in-depth:conformations}{chair} the methyl alone would cost about
\qty{5.7}{kJ/mol} and the larger isopropyl group more; the OH adds a smaller
cost of its own, and the \omterm{def:b1:stereochemistry-in-depth:conformations}{axial} OH and methyl, on carbons 1 and 3 and on the
same face, would also crowd each other. That is well over \qty{12}{kJ/mol},
a ratio of more than $\exp(12000/2479) \approx 130$ to one: menthol is
almost entirely in the \omterm{def:b1:stereochemistry-in-depth:conformations}{all-equatorial} \omterm{def:b1:stereochemistry-in-depth:conformations}{chair}.
\textbf{11.} A diastereomer: only one of three centres differs.
\textbf{12.} With isopropyl and methyl \omterm{def:b1:stereochemistry-in-depth:conformations}{equatorial}, the OH is \omterm{def:b1:stereochemistry-in-depth:conformations}{axial}.
\textbf{13.} \omterm{def:b1:stereochemistry-in-depth:enantiomers}{Diastereomers} differ in all their properties (here the
hydrogen bonding of an \omterm{def:b1:stereochemistry-in-depth:conformations}{axial} or \omterm{def:b1:stereochemistry-in-depth:conformations}{equatorial} OH); \omterm{def:b1:stereochemistry-in-depth:enantiomers}{enantiomers} differ only
towards \omterm{def:b1:stereochemistry-in-depth:stereocentre}{chiral} partners.
\textbf{14.} $[\alpha] = -2.46/(1.00 \times 0.0500) = -49.2^\circ$, inside the
handbook range.
% ledger: alpha:l-menthol-range
\textbf{15.} $+49.2^\circ$; $0$.
\textbf{16.} $-1.97/0.0500 = -39.4^\circ$.
\textbf{17.} $\alpha = \ell c[\alpha]_{(-)}\bigl(x - (1 - x)\bigr) =
\ell c[\alpha]_{(-)}(2x - 1)$.
\textbf{18.} $x = \frac12(1 + \alpha/\alpha_{\text{ref}})$ at equal $\ell$ and $c$.
\textbf{19.} Yes: any other optically active compound adds its own term
(\omterm{prop:b1:stereochemistry-in-depth:biot}{Biot's law}). A check of chemical purity by another method
(chromatography) is needed before reading the rotation as a composition.
\textbf{20.} $x = \frac12(1 + 1.97/2.46) = 0.900$.
\textbf{21.} $x = \frac12(1 + 2.40/2.46) = 0.988$.
\textbf{22.} The second (\qty{98.8}{\percent}).
\textbf{23.} $-3.94^\circ$; the same reading error is then a smaller
fraction of the angle.
\textbf{24.} \textbf{\qty{90.0}{\percent} of $(-)$-menthol} (and
\qty{10.0}{\percent} of its enantiomer).
\end{solution}
